\documentclass[11pt]{article}
\usepackage[a4paper,margin=25mm]{geometry}
\usepackage{amsmath,amssymb,amsthm}
\usepackage[hidelinks]{hyperref}
\emergencystretch=2em
\newtheorem{theorem}{Theorem}
\newtheorem{lemma}{Lemma}
\newcommand{\R}{\mathbb R}
\newcommand{\Z}{\mathbb Z}
\newcommand{\C}{\mathbb C}
\title{A uniform pole-free disk for cubic Ihara zeta functions\\
\large Disproof of Conjecture 00000007915}
\author{gaochengzhecpu}\date{}
\begin{document}\maketitle
\begin{abstract}
For every finite cubic simple graph, all Ihara poles have modulus at
least $1/2$. The circle proposed in the source has radius
$1/(2\sqrt2)<1/2$. A fixed positive gap separates every pole of every
such graph from that circle, ruling out the asserted equidistribution
even after only nontrivial poles are retained.
\end{abstract}

\section*{The asserted circle and the actual operator}
For $k=3$, the source proposes pole equidistribution on
\[
 S_r=\{u\in\C:|u|=r\},\qquad r=\frac1{2\sqrt2}.
\]
We use the unrescaled Ihara variable $u$ appearing in that assertion.
Let $D(G)$ be the oriented edges of a finite simple graph. Define its
nonbacktracking matrix by
\[
 B_{(a,b),(c,d)}=
 \begin{cases}1,&b=c\text{ and }d\ne a,\\0,&\text{otherwise.}\end{cases}
\]
If $G$ is cubic, each row has exactly two entries equal to $1$:
at the current endpoint there are three neighbors, of which immediate
reversal excludes one. In particular, every row sum is at most $2$.
The Hashimoto formula stated as the framework of the source is
\[
 \zeta_G(u)=\frac{1}{P_G(u)},\qquad P_G(u)=\det(I-uB).
\]
This identifies the operator we use with the genuine Ihara function.
Because the numerator is $1$ and $P_G(0)=1$, its finite poles are
exactly the zeros of $P_G$, counted with multiplicity. The determinant
formula is also recorded, for example, in\footnote{\href{https://www.kurims.kyoto-u.ac.jp/~kyodo/kokyuroku/contents/pdf/2120-10.pdf}{\emph{The Ihara Zeta function and Quantum Walk}, Theorem 3 (Hashimoto--Bass).}}
the standard graph-zeta literature.

\section*{A bound uniform over all cubic graphs}
\begin{lemma}
For every finite cubic simple graph $G$, $P_G(u)\ne0$ whenever
$|u|<1/2$.
\end{lemma}
\begin{proof}
If $P_G(u)=0$, there is a nonzero vector $x\in\C^{D(G)}$ such that
$(I-uB)x=0$, or $x=uBx$. Choose a coordinate $a$ for which
$M=|x_a|=\max_e|x_e|>0$. The row bound gives
\[
 M=|u|\,|(Bx)_a|
   \le |u|\sum_{e:B_{a,e}=1}|x_e|
   \le 2|u|M.
\]
Dividing by $M>0$ yields $|u|\ge1/2$, a contradiction.
\end{proof}

Set $\delta=1/2-r>0$. For every pole $u$ and every $w\in S_r$,
the reverse triangle inequality gives
\[
 |u-w|\ge |u|-|w|\ge\tfrac12-r=\delta.
\]
The same $\delta$ works for every graph size and every realization
of a random cubic graph. These bounds apply to all poles, and hence
also to any collection selected by nontrivial spectral parameters.

\section*{Why this disproves equidistribution}
This is a uniform obstruction, rather than an atypical finite example.
Define the bounded continuous function
\[
 \varphi(z)=\max\left\{0,\,1-
          \frac{\bigl||z|-r\bigr|}{\delta}\right\}.
\]
It equals $1$ at every point of $S_r$ and $0$ at every possible pole,
since those poles have modulus at least $1/2=r+\delta$.
Every normalized empirical pole measure therefore integrates
$\varphi$ to $0$, as does its expectation. Any probability measure
supported on $S_r$, in particular uniform angular measure, integrates
$\varphi$ to $1$. Weak convergence to such a measure is impossible.
The argument also rules out convergence in probability or almost
surely for random empirical measures: the test integral is identically
zero at every stage. Deleting trivial poles or changing multiplicities
cannot produce mass in the separated neighborhood of the circle.
Thus the first assertion of the source fails already at $k=3$.

\section*{Correspondence with Lean}
The matrix \texttt{nonbacktracking} is built on actual
\texttt{SimpleGraph.Dart} objects. The theorem
\texttt{successors\_\allowbreak{}card\_\allowbreak{}le} injects its legal successor set into
the current vertex's neighbor set with the reverse endpoint erased;
actual graph degree $3$ supplies the bound $2$.
The proof of \texttt{ihara\_\allowbreak{}denominator\_\allowbreak{}nonzero} uses Mathlib's
equivalence between zero matrix determinant and a nonzero kernel
vector, then the actual maximum-coordinate argument above.
The final theorem \texttt{uniform\_\allowbreak{}separation\_\allowbreak{}from\_\allowbreak{}claimed\_\allowbreak{}circle}
gives the positive gap for every genuine denominator root.
The complete graph $K_4$ formally witnesses that the cubic hypothesis
is realizable.

The formal result is the universal graph-determinant obstruction.
The Hashimoto identification of this denominator with Ihara zeta and
the bounded-continuous-test definition of weak convergence are the
explicit mathematical bridges supplied above. The Lean file does not
claim to formalize the general Ihara Euler-product theorem or a
random-graph limit theorem.

\section*{Reproduction}
The project pins Lean 4.19.0 and Mathlib commit
\\\texttt{c44e0c8ee63ca166450922a373c7409c5d26b00b}.\par
From \texttt{lean/}, run \texttt{lake build} and
\texttt{lake env lean -DwarningAsError=true Main.lean}.
On a new machine, \texttt{lake exe cache get} retrieves the official
pinned dependency artifacts. The submission itself is checked in a fresh
build directory. Printed dependencies use only the standard
\texttt{propext}, \texttt{Classical.choice}, and \texttt{Quot.sound}.
No additional axioms, proof placeholders, or numerical oracles are used.
Run \texttt{tectonic main.tex} to reproduce the PDF; no auxiliary numerical
script is needed. Build logs, source hashes, and visual-review records
are included under \texttt{verification/}.
\section*{Conjecture source}
\href{https://github.com/The-Last-Math-Competition/The-Last-Math-Competition/blob/main/conjectures/00000007915.md}{The Last Math Competition, Conjecture 00000007915}. The exact bilingual source and its commit and SHA-256 are preserved in this submission.
\end{document}
